\documentclass[border=10pt]{standalone}
\usepackage[T1]{fontenc}
\usepackage{lmodern}
\usepackage{tikz}
\usepackage{amsmath,amssymb}
\usetikzlibrary{calc}

\begin{document}
\begin{tikzpicture}

% Title
\node[font=\large\bfseries] at (5, 10.5) {Grid Cell Hexagonal Firing Pattern};

% --- Panel A: Spatial firing field ---
\node[font=\bfseries] at (2.5, 9.5) {(A) Spatial firing fields};

% Draw arena boundary
\draw[thick, gray] (0, 4.5) rectangle (5, 9);

% Hexagonal grid of firing fields
% Row spacing: dy = sqrt(3)/2 * spacing
% Using spacing = 1.2
\def\spacing{1.2}
\def\dy{1.039}  % sqrt(3)/2 * 1.2

% Generate hexagonal grid points
\foreach \row in {0,...,3} {
  \pgfmathsetmacro{\ypos}{5.0 + \row * \dy}
  \pgfmathsetmacro{\xoffset}{mod(\row,2) * \spacing/2}
  \foreach \col in {0,...,3} {
    \pgfmathsetmacro{\xpos}{0.8 + \col * \spacing + \xoffset}
    \pgfmathparse{\xpos < 4.8 && \ypos < 8.8 ? 1 : 0}
    \ifnum\pgfmathresult=1
      % Gaussian-like firing field (filled circle with gradient)
      \fill[red!60, opacity=0.7] (\xpos, \ypos) circle (0.3);
      \fill[red!90, opacity=0.9] (\xpos, \ypos) circle (0.15);
    \fi
  }
}

% Animal trajectory (light gray path)
\draw[gray!40, thin] (0.5, 5.0) -- (1.2, 5.8) -- (2.0, 5.3) -- (2.8, 6.2) -- 
  (3.5, 5.7) -- (4.2, 6.5) -- (3.8, 7.3) -- (3.0, 7.8) -- (2.2, 7.2) -- 
  (1.5, 7.8) -- (1.0, 7.0) -- (1.8, 6.5) -- (2.5, 7.0) -- (3.2, 6.8) --
  (4.0, 7.5) -- (4.5, 8.2) -- (3.8, 8.5) -- (3.0, 8.0) -- (2.0, 8.5) --
  (1.2, 8.0) -- (0.5, 8.5);

% Axis labels
\node[font=\small] at (2.5, 4.1) {Position in arena};

% --- Panel B: Hexagonal lattice with vectors ---
\node[font=\bfseries] at (9, 9.5) {(B) Hexagonal lattice structure};

% Draw hexagonal lattice with connecting lines
\def\hspacing{1.0}
\def\hdy{0.866}  % sqrt(3)/2

\foreach \row in {0,...,4} {
  \pgfmathsetmacro{\ypos}{4.8 + \row * \hdy}
  \pgfmathsetmacro{\xoffset}{mod(\row,2) * \hspacing/2}
  \foreach \col in {0,...,4} {
    \pgfmathsetmacro{\xpos}{7.0 + \col * \hspacing + \xoffset}
    \pgfmathparse{\xpos < 11.5 && \ypos < 9.2 ? 1 : 0}
    \ifnum\pgfmathresult=1
      \fill[blue!70] (\xpos, \ypos) circle (0.08);
    \fi
  }
}

% Draw lattice vectors from a central point
\coordinate (origin) at (8.5, 6.5);
\fill[black] (origin) circle (0.1);

% Vector e1
\draw[-{Stealth[length=3mm]}, thick, red!70!black] (origin) -- ++(1.0, 0) 
  node[below, font=\small] {$\mathbf{e}_1$};

% Vector e2
\draw[-{Stealth[length=3mm]}, thick, blue!70!black] (origin) -- ++(0.5, 0.866) 
  node[above left, font=\small] {$\mathbf{e}_2$};

% Angle arc
\draw[thick, gray] (8.9, 6.5) arc (0:60:0.4);
\node[font=\scriptsize, gray] at (9.15, 6.75) {$60^{\circ}$};

% Grid spacing annotation
\draw[<->, thick, gray!60] (7.0, 4.4) -- (8.0, 4.4);
\node[font=\small, gray!80] at (7.5, 4.0) {Grid spacing $\lambda$};

% --- Panel C: Phase wrapping on torus ---
\node[font=\bfseries] at (5, 3.5) {(C) Phase wrapping produces toroidal topology};

% Phase variable 1
\draw[-{Stealth[length=2.5mm]}, thick] (1.0, 2.5) -- (4.5, 2.5);
\node[font=\small] at (2.75, 2.1) {Spatial position $x$};

% Periodic wrapping marks
\foreach \x in {1.5, 2.7, 3.9} {
  \draw[red!70, thick, dashed] (\x, 2.3) -- (\x, 2.7);
}
\node[font=\scriptsize, red!70!black] at (2.1, 3.0) {$\lambda$};
\node[font=\scriptsize, red!70!black] at (3.3, 3.0) {$\lambda$};

% Arrow to circle
\draw[-{Stealth[length=2.5mm]}, thick, gray] (4.8, 2.5) -- (5.5, 2.5);

% Circle S^1
\draw[thick, blue!70] (6.3, 2.5) circle (0.6);
\draw[-{Stealth[length=2mm]}, red!70!black, thick] (6.3, 2.5) -- ++(60:0.6);
\node[font=\small] at (6.3, 1.5) {$\phi_1 \in S^1$};
\node[font=\scriptsize] at (5.3, 2.0) {mod $\lambda$};

% Cross product
\node[font=\Large] at (7.5, 2.5) {$\times$};

% Circle S^1 (second)
\draw[thick, green!60!black] (8.5, 2.5) circle (0.6);
\draw[-{Stealth[length=2mm]}, red!70!black, thick] (8.5, 2.5) -- ++(120:0.6);
\node[font=\small] at (8.5, 1.5) {$\phi_2 \in S^1$};

% Equals
\node[font=\Large] at (9.6, 2.5) {$=$};

% Torus T^2
\draw[thick, purple!70] (10.5, 2.5) ellipse (0.7 and 0.35);
\draw[thick, purple!70] (10.15, 2.5) arc (180:360:0.35 and 0.15);
\draw[thick, purple!70, dashed] (10.15, 2.5) arc (180:0:0.35 and 0.15);
\node[font=\small] at (10.5, 1.5) {$T^2 = S^1 \times S^1$};

% Bottom annotation
\node[font=\small, text width=9cm, align=center] at (5.5, 0.5) 
  {Periodic spatial coding in two directions produces two cyclic phase variables.\\
   Their joint state space is the two-torus $T^2$.};

\end{tikzpicture}
\end{document}
